% drafts/06-validation.tex — Empirical Validation
%
% Target length: ~2 pages (two-column PRA)
% Subsections: 6.1 Eigenspectrum, 6.2 Anticommutation,
%              6.3 Scaling, 6.4 Monotonicity census,
%              6.5 Parity operator

All computations are performed using the open-source \textsc{FockMap}
library~\cite{fockmap2026}.  No matrix representations are needed:
the Pauli algebra is computed symbolically via the \texttt{PauliRegister\-Sequence}
type, which implements exact finite-group multiplication.

%%====================================================================
\subsection{Eigenspectrum equivalence}
\label{sec:val-eigen}

We encode the H$_2$ STO-3G Hamiltonian (4~spin-orbitals, nuclear
separation $R = 0.7414\;\text{\AA}$) with all five encodings.  Each
produces a 15-term qubit Hamiltonian with identical summed absolute
coefficients $\sum |c_i| = 3.2557$.

Full $16 \times 16$ matrix diagonalisation confirms that all five
encodings yield identical eigenspectra to machine precision
($|\Delta\lambda| = 4.44 \times 10^{-16}$).  The ground-state energy
is $E_0 = -1.7622\;\text{Ha}$ (electronic FCI), giving
$E_\text{total} = -1.0471\;\text{Ha}$ with nuclear repulsion.

The Jordan--Wigner encoding produces the expected excitation terms
$\pm 0.1744\; XXYY$, $\pm 0.1744\; YYXX$, while the ternary tree
encoding distributes these among different Pauli signatures ($XXXY$,
$XXYX$, $XYXX$, $XYYY$) — same physics, different representation.

%%====================================================================
\subsection{Anticommutation verification}
\label{sec:val-car}

For each encoding, we construct the full set of $2n$ Majorana operators
and verify all $\binom{2n}{2}$ anticommutators both symbolically (via
\textsc{FockMap}'s Pauli algebra) and by matrix computation at
small~$n$.

\begin{itemize}
  \item $n = 4, 8, 16$: all five path-based encodings satisfy the CAR
    with zero deviation.
  \item $n = 8$, balanced ternary tree with Construction~A:
    $\acomm{\an{4}}{\ad{7}} \neq 0$, confirming the star-tree
    restriction (\cref{thm:monotonicity}).
  \item $n = 8$, Fenwick tree with Construction~A (generic
    \texttt{treeEncodingScheme}, \emph{not} the BK-specific formulas):
    CAR violated --- confirming that Construction~F is distinct from
    Construction~A applied to a Fenwick-shaped tree.
\end{itemize}

%%====================================================================
\subsection{Scaling benchmark}
\label{sec:val-scaling}

\Cref{tab:scaling} reports the maximum Pauli weight of any single
ladder operator for $n$ up to~24.

\begin{table}[t]
\centering
\caption{Maximum Pauli weight $\max_j \pauliw{\ad{j}}$ across encodings.}
\label{tab:scaling}
\begin{tabular}{@{}rrrrrr@{}}
\toprule
$n$ & JW & BK & Parity & BinTree & TerTree \\
\midrule
 4  &  4 &  3 &  4 &  3 & \textbf{2} \\
 8  &  8 &  4 &  8 &  4 & \textbf{3} \\
12  & 12 &  4 & 12 &  4 & \textbf{4} \\
16  & 16 &  5 & 16 &  5 & \textbf{4} \\
20  & 20 &  5 & 20 &  5 & \textbf{4} \\
24  & 24 &  5 & 24 &  5 & \textbf{5} \\
\bottomrule
\end{tabular}
\end{table}

A log-log fit confirms slopes consistent with $O(n)$ for JW/Parity,
$O(\log_2 n)$ for BK/BinTree, and $O(\log_3 n)$ for TerTree.

\Cref{tab:mean-weight} reports mean creation-operator weights,
showing that the ternary tree has the lowest average weight at every
system size.

\begin{table}[t]
\centering
\caption{Mean Pauli weight $\frac{1}{n}\sum_j \pauliw{\ad{j}}$ across encodings.}
\label{tab:mean-weight}
\begin{tabular}{@{}rrrrrr@{}}
\toprule
$n$ & JW & BK & Parity & BinTree & TerTree \\
\midrule
 4  & 2.50 & 2.62 & 2.88 & 2.38 & \textbf{2.00} \\
 8  & 4.50 & 3.56 & 4.94 & 3.06 & \textbf{2.75} \\
16  & 8.50 & 4.53 & 8.97 & 3.84 & \textbf{3.28} \\
24  & 12.50& 4.54 &12.98 & 4.40 & \textbf{3.71} \\
\bottomrule
\end{tabular}
\end{table}

% TODO: Figure — log-log plot with fitted slopes.

%%====================================================================
\subsection{Parity operator weights}
\label{sec:val-parity}

Using symbolic Pauli algebra (no matrices), we compute the encoded
parity operator $\hat{P} = E(\prod_j (I - 2\nhat{j}))$ for each
encoding.  All parity operators are single Pauli strings ($P^2 = I$
verified symbolically).

\begin{table}[t]
\centering
\caption{Parity operator $\hat{P}$ weight and structure.}
\label{tab:parity}
\begin{tabular}{@{}lrl@{}}
\toprule
\textbf{Encoding} & $\pauliw{\hat{P}}$ & \textbf{Structure} \\
\midrule
Jordan--Wigner     & $n$            & $Z^{\otimes n}$ (all-$Z$ string) \\
Parity             & $1$            & $Z_{n-1}$ (single qubit) \\
Balanced binary    & $1$            & $Z_\text{root}$ \\
Bravyi--Kitaev     & $O(\log n)$    & $Z$ on Fenwick path to root \\
Balanced ternary   & $O(\log_3 n)$  & $Z$ on ternary root path \\
\bottomrule
\end{tabular}
\end{table}

This confirms Theorem~\ref{thm:parity-weight}: the parity weight is
determined by the $Z$-chain structure of the tree.  The computation
scales to $n = 100$ in under 6~seconds using purely symbolic algebra
--- demonstrating the library's practical utility for large systems.

%%====================================================================
\subsection{Monotonicity census and star-tree enumeration}
\label{sec:val-census}

We enumerate all $n^{n-1}$ labelled rooted trees for $n = 1, \ldots, 6$
via Pr\"ufer sequences, classify each as monotonic or non-monotonic,
and test Construction~A against the full CAR.

\begin{table}[t]
\centering
\caption{Monotonicity census and Construction~A domain.  $|M(n)|$ =
  monotonic trees; ``Pass CAR'' = trees where Construction~A satisfies
  all anticommutation relations.}
\label{tab:monotonicity}
\begin{tabular}{@{}rrrrr@{}}
\toprule
$n$ & $n^{n-1}$ & $|M(n)|$ & $(n{-}1)!$ & Pass CAR \\
\midrule
1 &       1 &     1 &     1 &  1 \\
2 &       2 &     1 &     1 &  2 \\
3 &       9 &     2 &     2 &  3 \\
4 &      64 &     6 &     6 &  4 \\
5 &     625 &    24 &    24 &  5 \\
6 &    7776 &   120 &   120 &  6 \\
\bottomrule
\end{tabular}
\end{table}

Two results are visible in~\cref{tab:monotonicity}:
\begin{enumerate}
  \item $|M(n)| = (n{-}1)!$ exactly, for every $n$ tested
    (\cref{prop:monotonic-count}).
  \item The number of trees passing the CAR under Construction~A
    is exactly~$n$ --- the $n$ possible star trees.  This is strictly
    less than $(n{-}1)!$ for $n \ge 4$, confirming that monotonicity
    is necessary but not sufficient (\cref{thm:monotonicity}).
\end{enumerate}
